\Askhsh{14970}{4}
\begin{ylh}{1}
    \item 2.1. Εξίσωση Ευθείας.
\end{ylh}
Δίνεται η ευθεία $y=\lambda(x-2)+\lambda-2, \lambda\in\mathbb{R}$ (1).

\begin{alist}
\item Να βρείτε τις ευθείες που προκύπτουν όταν $\lambda=1$ και όταν $\lambda=2$. Κατόπιν να βρείτε το κοινό σημείο $M$ των δύο ευθειών.  \monades{7}

Έστω $M(1, -2)$.

\item Να αποδείξετε ότι όλες οι ευθείες που προκύπτουν από την (1) για τις διάφορες τιμές του $\lambda$, διέρχονται από το $M$.  \monades{5}

\item Να βρείτε:
\begin{rlist}
\item τα σημεία τομής $A, B$ της ευθείας (1) με τους άξονες $x'x$ και $y'y$.  \monades{6}
\item για ποιες τιμές του $\lambda$ το εμβαδόν του τριγώνου $OAB$ είναι ίσο με $\frac{1}{2}$.  \monades{7}
\end{rlist}
\end{alist}